% --- Archon physics formalization source begin ---
% archon:physics
% archon:covers PhyXMiniProblems/problem_phyx_mini_0122.lean
% archon:source-report reports/phyx_mini/problem_phyx_mini_0122.source.json

\chapter{Physics problem phyx\_mini\_0122}
\label{ch:PhyXMiniProblems_problem_phyx_mini_0122}

\paragraph{Problem source.}
\#\# Physical scenario

While researching the use of laser pointers, you conduct a diffraction experiment with two thin parallel slits. Your result is the pattern of closely spaced bright and dark fringes shown in figure. (Only the central portion of the pattern is shown.) You measure that the bright spots are equally spaced at 1.53 cm center to center (except for the missing spots) on a screen that is 2.50 m from the slits. The light source was a helium-neon laser producing a wavelength of 632.8 nm.

\#\# Question

How wide is each one?

\#\# Answer choices

- A: A: 0.148mm

- B: B: 0.152mm

- C: C: 0.134mm

- D: D: 0.132mm

\#\# Figure caption (auxiliary; use the image as primary evidence)

The image shows a black rectangular background with a horizontal row of evenly spaced white dots. Above the center of the row, text reads "1.53 mm," accompanied by two arrows pointing towards each other, indicating the distance between the centers of two adjacent dots. The illustration measures the spacing between the dots.

\#\# Dataset metadata

Wave Optics; Physical Model Grounding Reasoning, Multi-Formula Reasoning

\paragraph{Recorded answer/context.}
A: A: 0.148mm

\paragraph{Figure/image path.}
/root/proposal\_for\_physic/science-mango/phyx\_mini\_run/phyx\_data/test\_image/122.png

\paragraph{Formalization target.}
create a compiling Lean file with sorry bodies at `PhyXMiniProblems/problem\_phyx\_mini\_0122.lean`. The Lean declarations must preserve the physical quantities, dimensions or dimensional roles, figure labels, governing-law hypotheses, and final relation expressed by this problem.
Use Mathlib/Physlib names found through LeanExplore where available. If a physics API is missing, introduce faithful local abstractions rather than scalar placeholder aliases.

\begin{theorem}[Physics formalization target]
\label{thm:physics:phyx_mini_0122:target}
\lean{PhyXMiniProblems.ProblemPhyXMini0122.problem_phyx_mini_0122}
\uses{def:physics:phyx-mini-0122:phyxminiproblems-problemphyxmini0122-lengthinmillimeters, def:physics:phyx-mini-0122:phyxminiproblems-problemphyxmini0122-missingfringedoubleslitsetup, def:physics:phyx-mini-0122:phyxminiproblems-problemphyxmini0122-hasphysicalparameters, def:physics:phyx-mini-0122:phyxminiproblems-problemphyxmini0122-matchesslitgeometry, def:physics:phyx-mini-0122:phyxminiproblems-problemphyxmini0122-matchesproblemreadouts, def:physics:phyx-mini-0122:phyxminiproblems-problemphyxmini0122-matchesfigurereadouts, def:physics:phyx-mini-0122:phyxminiproblems-problemphyxmini0122-satisfiesconstructiveinterferencelaw, def:physics:phyx-mini-0122:phyxminiproblems-problemphyxmini0122-satisfiesplanarscreenprojection, def:physics:phyx-mini-0122:phyxminiproblems-problemphyxmini0122-satisfiesfirstminimumdiffractionlaw, def:physics:phyx-mini-0122:phyxminiproblems-problemphyxmini0122-missingspotscoincidewithfirstminima, def:physics:phyx-mini-0122:phyxminiproblems-problemphyxmini0122-isuniquematchingdisplayedwidth}
The assigned autoformalize agent should translate this physics problem into Lean declarations in the covered file, with theorem and lemma proof bodies written as `by sorry`.
\end{theorem}
\begin{proof}
This is an autoformalization task, not a proof task. Produce statements that can later be proved by the physics prover without weakening the physical model.
\end{proof}

% --- Archon declaration topology begin ---
\section{Declaration topology}
\begin{definition}[Length Quantity]
  \label{def:physics:phyx-mini-0122:phyxminiproblems-problemphyxmini0122-lengthquantity}
  \lean{PhyXMiniProblems.ProblemPhyXMini0122.LengthQuantity}
  A signed, unit-independent physical quantity carrying length dimension
\end{definition}
\begin{proof}
  This is the declared model definition of Length Quantity.
\end{proof}
\begin{definition}[length Readout]
  \label{def:physics:phyx-mini-0122:phyxminiproblems-problemphyxmini0122-lengthreadout}
  \lean{PhyXMiniProblems.ProblemPhyXMini0122.lengthReadout}
  \uses{def:physics:phyx-mini-0122:phyxminiproblems-problemphyxmini0122-lengthquantity}
  Read a physical length as a real number in the selected unit
\end{definition}
\begin{proof}
  This is the declared model definition of length Readout.
\end{proof}
\begin{definition}[length In Meters]
  \label{def:physics:phyx-mini-0122:phyxminiproblems-problemphyxmini0122-lengthinmeters}
  \lean{PhyXMiniProblems.ProblemPhyXMini0122.lengthInMeters}
  \uses{def:physics:phyx-mini-0122:phyxminiproblems-problemphyxmini0122-lengthquantity, def:physics:phyx-mini-0122:phyxminiproblems-problemphyxmini0122-lengthreadout}
  The SI-metre readout of a physical length
\end{definition}
\begin{proof}
  This is the declared model definition of length In Meters.
\end{proof}
\begin{definition}[length In Centimeters]
  \label{def:physics:phyx-mini-0122:phyxminiproblems-problemphyxmini0122-lengthincentimeters}
  \lean{PhyXMiniProblems.ProblemPhyXMini0122.lengthInCentimeters}
  \uses{def:physics:phyx-mini-0122:phyxminiproblems-problemphyxmini0122-lengthquantity, def:physics:phyx-mini-0122:phyxminiproblems-problemphyxmini0122-lengthreadout}
  The centimetre readout of a physical length
\end{definition}
\begin{proof}
  This is the declared model definition of length In Centimeters.
\end{proof}
\begin{definition}[length In Millimeters]
  \label{def:physics:phyx-mini-0122:phyxminiproblems-problemphyxmini0122-lengthinmillimeters}
  \lean{PhyXMiniProblems.ProblemPhyXMini0122.lengthInMillimeters}
  \uses{def:physics:phyx-mini-0122:phyxminiproblems-problemphyxmini0122-lengthquantity, def:physics:phyx-mini-0122:phyxminiproblems-problemphyxmini0122-lengthreadout}
  The millimetre readout of a physical length
\end{definition}
\begin{proof}
  This is the declared model definition of length In Millimeters.
\end{proof}
\begin{definition}[length In Nanometers]
  \label{def:physics:phyx-mini-0122:phyxminiproblems-problemphyxmini0122-lengthinnanometers}
  \lean{PhyXMiniProblems.ProblemPhyXMini0122.lengthInNanometers}
  \uses{def:physics:phyx-mini-0122:phyxminiproblems-problemphyxmini0122-lengthquantity, def:physics:phyx-mini-0122:phyxminiproblems-problemphyxmini0122-lengthreadout}
  The nanometre readout of a physical length
\end{definition}
\begin{proof}
  This is the declared model definition of length In Nanometers.
\end{proof}
\begin{definition}[Slit Label]
  \label{def:physics:phyx-mini-0122:phyxminiproblems-problemphyxmini0122-slitlabel}
  \lean{PhyXMiniProblems.ProblemPhyXMini0122.SlitLabel}
  Labels for the two parallel openings
\end{definition}
\begin{proof}
  This is the declared model definition of Slit Label.
\end{proof}
\begin{definition}[Screen Side]
  \label{def:physics:phyx-mini-0122:phyxminiproblems-problemphyxmini0122-screenside}
  \lean{PhyXMiniProblems.ProblemPhyXMini0122.ScreenSide}
  The two sides of the symmetric screen pattern
\end{definition}
\begin{proof}
  This is the declared model definition of Screen Side.
\end{proof}
\begin{definition}[Spot Visibility]
  \label{def:physics:phyx-mini-0122:phyxminiproblems-problemphyxmini0122-spotvisibility}
  \lean{PhyXMiniProblems.ProblemPhyXMini0122.SpotVisibility}
  Whether an ideal constructive order is visible in the cropped image
\end{definition}
\begin{proof}
  This is the declared model definition of Spot Visibility.
\end{proof}
\begin{definition}[Slit Geometry]
  \label{def:physics:phyx-mini-0122:phyxminiproblems-problemphyxmini0122-slitgeometry}
  \lean{PhyXMiniProblems.ProblemPhyXMini0122.SlitGeometry}
  Qualitative geometry assigned to the two openings in the experiment
\end{definition}
\begin{proof}
  This is the declared model definition of Slit Geometry.
\end{proof}
\begin{definition}[Screen Geometry]
  \label{def:physics:phyx-mini-0122:phyxminiproblems-problemphyxmini0122-screengeometry}
  \lean{PhyXMiniProblems.ProblemPhyXMini0122.ScreenGeometry}
  Geometry used to project an outgoing direction onto the screen
\end{definition}
\begin{proof}
  This is the declared model definition of Screen Geometry.
\end{proof}
\begin{definition}[Spacing Evidence Policy]
  \label{def:physics:phyx-mini-0122:phyxminiproblems-problemphyxmini0122-spacingevidencepolicy}
  \lean{PhyXMiniProblems.ProblemPhyXMini0122.SpacingEvidencePolicy}
  Which of the inconsistent spacing reports controls the optical model
\end{definition}
\begin{proof}
  This is the declared model definition of Spacing Evidence Policy.
\end{proof}
\begin{definition}[Missing Fringe Double Slit Setup]
  \label{def:physics:phyx-mini-0122:phyxminiproblems-problemphyxmini0122-missingfringedoubleslitsetup}
  \lean{PhyXMiniProblems.ProblemPhyXMini0122.MissingFringeDoubleSlitSetup}
  \uses{def:physics:phyx-mini-0122:phyxminiproblems-problemphyxmini0122-lengthquantity, def:physics:phyx-mini-0122:phyxminiproblems-problemphyxmini0122-slitlabel, def:physics:phyx-mini-0122:phyxminiproblems-problemphyxmini0122-screenside, def:physics:phyx-mini-0122:phyxminiproblems-problemphyxmini0122-spotvisibility, def:physics:phyx-mini-0122:phyxminiproblems-problemphyxmini0122-slitgeometry, def:physics:phyx-mini-0122:phyxminiproblems-problemphyxmini0122-screengeometry, def:physics:phyx-mini-0122:phyxminiproblems-problemphyxmini0122-spacingevidencepolicy}
  Physical quantities, order-indexed angles and screen locations, and figure metadata for the experiment. screenOffset order is the signed position of the ideal constructive maximum of that order, even when diffraction suppresses the spot. slitWidth is an unknown physical length and no field assigns its requested numerical value
\end{definition}
\begin{proof}
  This is the declared model definition of Missing Fringe Double Slit Setup.
\end{proof}
\begin{definition}[Has Physical Parameters]
  \label{def:physics:phyx-mini-0122:phyxminiproblems-problemphyxmini0122-hasphysicalparameters}
  \lean{PhyXMiniProblems.ProblemPhyXMini0122.HasPhysicalParameters}
  \uses{def:physics:phyx-mini-0122:phyxminiproblems-problemphyxmini0122-lengthinmeters, def:physics:phyx-mini-0122:phyxminiproblems-problemphyxmini0122-missingfringedoubleslitsetup}
  Positive physical lengths and the signs selecting the left and right angular branches. These conditions do not determine the unknown width
\end{definition}
\begin{proof}
  This is the declared model definition of Has Physical Parameters.
\end{proof}
\begin{definition}[Matches Slit Geometry]
  \label{def:physics:phyx-mini-0122:phyxminiproblems-problemphyxmini0122-matchesslitgeometry}
  \lean{PhyXMiniProblems.ProblemPhyXMini0122.MatchesSlitGeometry}
  \uses{def:physics:phyx-mini-0122:phyxminiproblems-problemphyxmini0122-lengthreadout, def:physics:phyx-mini-0122:phyxminiproblems-problemphyxmini0122-missingfringedoubleslitsetup}
  The openings are interpreted as parallel and equal-width, and the difference of their transverse center coordinates is their center separation. The common width remains unconstrained numerically
\end{definition}
\begin{proof}
  This is the declared model definition of Matches Slit Geometry.
\end{proof}
\begin{definition}[Matches Problem Readouts]
  \label{def:physics:phyx-mini-0122:phyxminiproblems-problemphyxmini0122-matchesproblemreadouts}
  \lean{PhyXMiniProblems.ProblemPhyXMini0122.MatchesProblemReadouts}
  \uses{def:physics:phyx-mini-0122:phyxminiproblems-problemphyxmini0122-lengthinmeters, def:physics:phyx-mini-0122:phyxminiproblems-problemphyxmini0122-lengthincentimeters, def:physics:phyx-mini-0122:phyxminiproblems-problemphyxmini0122-lengthinnanometers, def:physics:phyx-mini-0122:phyxminiproblems-problemphyxmini0122-missingfringedoubleslitsetup}
  Problem-text readouts: a helium-neon wavelength of 632.8 nm, a screen distance of 2.50 m, and the prose's adjacent-fringe report of 1.53 cm. The prose spacing is retained for provenance but is not the primary spacing used by the laws
\end{definition}
\begin{proof}
  This is the declared model definition of Matches Problem Readouts.
\end{proof}
\begin{definition}[Matches Figure Readouts]
  \label{def:physics:phyx-mini-0122:phyxminiproblems-problemphyxmini0122-matchesfigurereadouts}
  \lean{PhyXMiniProblems.ProblemPhyXMini0122.MatchesFigureReadouts}
  \uses{def:physics:phyx-mini-0122:phyxminiproblems-problemphyxmini0122-lengthinmillimeters, def:physics:phyx-mini-0122:phyxminiproblems-problemphyxmini0122-missingfringedoubleslitsetup}
  Figure readouts and labels. The cropped row represents ideal orders -9 through 9; the spots at -7 and 7 are missing, all other orders in the window are bright, and the inward arrows between orders -1 and 0 label a center-to-center spacing of 1.53 mm
\end{definition}
\begin{proof}
  This is the declared model definition of Matches Figure Readouts.
\end{proof}
\begin{definition}[Satisfies Constructive Interference Law]
  \label{def:physics:phyx-mini-0122:phyxminiproblems-problemphyxmini0122-satisfiesconstructiveinterferencelaw}
  \lean{PhyXMiniProblems.ProblemPhyXMini0122.SatisfiesConstructiveInterferenceLaw}
  \uses{def:physics:phyx-mini-0122:phyxminiproblems-problemphyxmini0122-lengthreadout, def:physics:phyx-mini-0122:phyxminiproblems-problemphyxmini0122-missingfringedoubleslitsetup}
  Constructive two-slit interference law d sin θₘ = m λ. It is stated in every length unit for the finite order window represented by the cropped figure. Restricting the law to physically represented orders is essential: no fixed positive d and λ could realize every integer order. The law does not contain the requested slit width
\end{definition}
\begin{proof}
  This is the declared model definition of Satisfies Constructive Interference Law.
\end{proof}
\begin{definition}[Satisfies Planar Screen Projection]
  \label{def:physics:phyx-mini-0122:phyxminiproblems-problemphyxmini0122-satisfiesplanarscreenprojection}
  \lean{PhyXMiniProblems.ProblemPhyXMini0122.SatisfiesPlanarScreenProjection}
  \uses{def:physics:phyx-mini-0122:phyxminiproblems-problemphyxmini0122-lengthreadout, def:physics:phyx-mini-0122:phyxminiproblems-problemphyxmini0122-missingfringedoubleslitsetup}
  Exact planar-screen projection yₘ = L tan θₘ. The screen is perpendicular to the central optical axis, so elementary ray geometry gives this tangent relation without a small-angle replacement. The ideal order position exists even when the finite slit width suppresses its brightness. Together with adjacent order positions, this law calibrates the slit separation but says nothing directly about slit width
\end{definition}
\begin{proof}
  This is the declared model definition of Satisfies Planar Screen Projection.
\end{proof}
\begin{definition}[Satisfies First Minimum Diffraction Law]
  \label{def:physics:phyx-mini-0122:phyxminiproblems-problemphyxmini0122-satisfiesfirstminimumdiffractionlaw}
  \lean{PhyXMiniProblems.ProblemPhyXMini0122.SatisfiesFirstMinimumDiffractionLaw}
  \uses{def:physics:phyx-mini-0122:phyxminiproblems-problemphyxmini0122-lengthreadout, def:physics:phyx-mini-0122:phyxminiproblems-problemphyxmini0122-screenside, def:physics:phyx-mini-0122:phyxminiproblems-problemphyxmini0122-missingfringedoubleslitsetup}
  First-minimum law for the diffraction envelope of each finite slit: a |sin θ₁| = λ. This is a governing single-slit law, not the requested numerical value of a
\end{definition}
\begin{proof}
  This is the declared model definition of Satisfies First Minimum Diffraction Law.
\end{proof}
\begin{definition}[Missing Spots Coincide With First Minima]
  \label{def:physics:phyx-mini-0122:phyxminiproblems-problemphyxmini0122-missingspotscoincidewithfirstminima}
  \lean{PhyXMiniProblems.ProblemPhyXMini0122.MissingSpotsCoincideWithFirstMinima}
  \uses{def:physics:phyx-mini-0122:phyxminiproblems-problemphyxmini0122-missingfringedoubleslitsetup}
  Physical interpretation of the missing spots: the first missing constructive order on each side coincides with the first minimum of the one-slit envelope. This identifies the relevant orders and angles but does not assume a width
\end{definition}
\begin{proof}
  This is the declared model definition of Missing Spots Coincide With First Minima.
\end{proof}
\begin{definition}[Answer Choice]
  \label{def:physics:phyx-mini-0122:phyxminiproblems-problemphyxmini0122-answerchoice}
  \lean{PhyXMiniProblems.ProblemPhyXMini0122.AnswerChoice}
  Labels of the four slit-width choices printed with the dataset item
\end{definition}
\begin{proof}
  This is the declared model definition of Answer Choice.
\end{proof}
\begin{definition}[displayed Width In Millimeters]
  \label{def:physics:phyx-mini-0122:phyxminiproblems-problemphyxmini0122-displayedwidthinmillimeters}
  \lean{PhyXMiniProblems.ProblemPhyXMini0122.displayedWidthInMillimeters}
  \uses{def:physics:phyx-mini-0122:phyxminiproblems-problemphyxmini0122-answerchoice}
  The width in millimetres printed beside each answer label
\end{definition}
\begin{proof}
  This is the declared model definition of displayed Width In Millimeters.
\end{proof}
\begin{definition}[recorded Dataset Answer]
  \label{def:physics:phyx-mini-0122:phyxminiproblems-problemphyxmini0122-recordeddatasetanswer}
  \lean{PhyXMiniProblems.ProblemPhyXMini0122.recordedDatasetAnswer}
  \uses{def:physics:phyx-mini-0122:phyxminiproblems-problemphyxmini0122-answerchoice}
  The answer label recorded by the source dataset
\end{definition}
\begin{proof}
  This is the declared model definition of recorded Dataset Answer.
\end{proof}
\begin{definition}[rounded To Nearest Thousandth]
  \label{def:physics:phyx-mini-0122:phyxminiproblems-problemphyxmini0122-roundedtonearestthousandth}
  \lean{PhyXMiniProblems.ProblemPhyXMini0122.roundedToNearestThousandth}
  Round a millimetre readout to the nearest thousandth of a millimetre
\end{definition}
\begin{proof}
  This is the declared model definition of rounded To Nearest Thousandth.
\end{proof}
\begin{definition}[Matches Displayed Width]
  \label{def:physics:phyx-mini-0122:phyxminiproblems-problemphyxmini0122-matchesdisplayedwidth}
  \lean{PhyXMiniProblems.ProblemPhyXMini0122.MatchesDisplayedWidth}
  \uses{def:physics:phyx-mini-0122:phyxminiproblems-problemphyxmini0122-lengthinmillimeters, def:physics:phyx-mini-0122:phyxminiproblems-problemphyxmini0122-missingfringedoubleslitsetup, def:physics:phyx-mini-0122:phyxminiproblems-problemphyxmini0122-answerchoice, def:physics:phyx-mini-0122:phyxminiproblems-problemphyxmini0122-displayedwidthinmillimeters, def:physics:phyx-mini-0122:phyxminiproblems-problemphyxmini0122-roundedtonearestthousandth}
  A displayed choice agrees with the derived rounded width
\end{definition}
\begin{proof}
  This is the declared model definition of Matches Displayed Width.
\end{proof}
\begin{definition}[Is Unique Matching Displayed Width]
  \label{def:physics:phyx-mini-0122:phyxminiproblems-problemphyxmini0122-isuniquematchingdisplayedwidth}
  \lean{PhyXMiniProblems.ProblemPhyXMini0122.IsUniqueMatchingDisplayedWidth}
  \uses{def:physics:phyx-mini-0122:phyxminiproblems-problemphyxmini0122-missingfringedoubleslitsetup, def:physics:phyx-mini-0122:phyxminiproblems-problemphyxmini0122-answerchoice, def:physics:phyx-mini-0122:phyxminiproblems-problemphyxmini0122-matchesdisplayedwidth}
  The selected answer is the unique displayed choice matching the width
\end{definition}
\begin{proof}
  This is the declared model definition of Is Unique Matching Displayed Width.
\end{proof}
\begin{lemma}[slit Width In Millimeters eq]
  \label{lem:physics:phyx-mini-0122:phyxminiproblems-problemphyxmini0122-slitwidthinmillimeters-eq}
  \lean{PhyXMiniProblems.ProblemPhyXMini0122.slitWidthInMillimeters_eq}
  \uses{def:physics:phyx-mini-0122:phyxminiproblems-problemphyxmini0122-lengthinmillimeters, def:physics:phyx-mini-0122:phyxminiproblems-problemphyxmini0122-missingfringedoubleslitsetup, def:physics:phyx-mini-0122:phyxminiproblems-problemphyxmini0122-hasphysicalparameters, def:physics:phyx-mini-0122:phyxminiproblems-problemphyxmini0122-matchesproblemreadouts, def:physics:phyx-mini-0122:phyxminiproblems-problemphyxmini0122-matchesfigurereadouts, def:physics:phyx-mini-0122:phyxminiproblems-problemphyxmini0122-satisfiesconstructiveinterferencelaw, def:physics:phyx-mini-0122:phyxminiproblems-problemphyxmini0122-satisfiesplanarscreenprojection, def:physics:phyx-mini-0122:phyxminiproblems-problemphyxmini0122-satisfiesfirstminimumdiffractionlaw, def:physics:phyx-mini-0122:phyxminiproblems-problemphyxmini0122-missingspotscoincidewithfirstminima}
  Using the figure-primary 1.53 mm spacing, λ = 632.8 nm, and L = 2.50 m, the interference and exact tangent-projection laws give the slit separation. Coincidence of order 7 with the first diffraction minimum then gives the exact common width (791/5355) * sqrt (1 + (153/250000)\textasciicircum{}2) mm, approximately 0.147712 mm
\end{lemma}
\begin{proof}
  The conclusion follows from the governing relations and figure readouts named in the statement of slit Width In Millimeters eq.
\end{proof}
% --- Archon declaration topology end ---
% --- Archon physics formalization source end ---
